\section{硬件亲和：算法必须能并行、能矩阵乘、能 scale}

前面讲了算法考古，本章进入能否落地的最后一道门槛：硬件。一个想法如果不能并行、不能写成高效 kernel、不能贴近现代 GPU 的矩阵乘能力，就很难进入 frontier 模型。

最后一条技术主线是硬件亲和。杨松琳认为，一个算法不仅要在 machine learning 上 make sense，还要能并行、能 scale、能在硬件上高效实现。Transformer 当年能成功，不只是因为表达能力强，也因为 attention 和 FFN 都有大量矩阵乘，硬件亲和优于 RNN/LSTM。

\begin{figure}[H]
\centering
\includegraphics[width=0.96\textwidth]{hardware-friendly-algorithm.png}
\caption{硬件亲和算法：数学 make sense，还要能并行、能矩阵乘、能 scale。自制概念图，依据 01:29:50--01:42:23 对谈内容整理。}
\end{figure}

\begin{knowledgebox}{读图：现代架构研究是算法和系统共同设计}
数学合理只是第一步；还要有并行算法，如 chunkwise；要能利用矩阵乘；要有 kernel 实现；最后还要在大模型规模上 scale。任何一环缺失，算法都难以进入 frontier 模型。
\end{knowledgebox}

\subsection{Chunkwise Algorithm 与并行}

本节把硬件亲和具体化为并行算法。线性注意力常有递推形式，递推本身不适合 GPU 大规模并行，所以必须把它改写成可分块并行的形式。

Linear Attention 常常有 recurrent 形式，但 recurrence 天然难并行。要让它在现代 GPU 上高效训练，就需要 chunkwise algorithm 等并行算法，把长序列切成块，既保留递推结构，又能利用并行计算。这个角色类似 FlashAttention 对 Softmax Attention 的作用：不是改变模型目标，而是让算法在硬件上跑得起来。

\begin{importantbox}{并行算法的意义}
没有高效并行算法，理论上很美的 attention 变体也无法进入大规模训练。架构创新和系统优化不是先后关系，而是共同决定一个想法能否落地。
\end{importantbox}

\subsection{DeepSeek 的硬件协同}

接下来用 DeepSeek 作为硬件协同样本。它的价值不只是选择 sparse，而是让 sparse 的实现形式贴近 FP8、矩阵乘和去 softmax 的硬件路径。

访谈中认为 DeepSeek 非常重视硬件和算法协同。DeepSeek Sparse Attention 的 indexer 可以用 FP8 计算 attention score / logit，去掉昂贵的 softmax/exp 操作，留下大量矩阵乘和 Top-K 选择。这使得它更贴近硬件原则，也更容易和 infra 团队协同。

\begin{figure}[H]
\centering
\includegraphics[width=0.94\textwidth]{deepseek-hardware-codesign.png}
\caption{DeepSeek 的硬件协同：稀疏选择、FP8 和矩阵乘共同服务效率。自制概念图，依据 01:29:50--01:42:23 对谈内容整理。}
\end{figure}

\begin{knowledgebox}{读图：硬件协同不是只写 kernel}
算法侧通过 Sparse / Indexer / Top-K 减少计算；硬件侧用 FP8、矩阵乘、去掉昂贵 softmax/exp 等方式提高效率。真正的 co-design 是算法目标和硬件原则同时考虑。
\end{knowledgebox}

\subsection{本章小结}

未来架构研究必须同时满足三件事：机器学习上合理，数学上可解释，硬件上能并行和矩阵乘。否则再漂亮的算法也很难进入 frontier 模型。

